% !TeX root = Book.Stress_regulation.tex
\documentclass[../Book.Stress_regulation.tex]{subfiles}
\graphicspath{{\subfix{../images/}}}
\begin{document}
	
\chapter{Emotions and Organs}

We recognize which organs are burdened worst by looking at the type of emotions, 
which reign that person.
If this emotion is anger, then liver and gall bladder are most affected. 
Pain will most likely occur in the middle of the back. 
If the patient has sorrows and is full of sadness, the lungs and large intestines are most affected by that.
The back pain sits higher up, immediately below the shoulders.
If there's excessive fear or melancholy, then the kidneys are affected.
This concerns the lower back, pain in the lower back is often the case.
Somebody that tends towards excessive sympathy for other and empathy and a weak personality,
there are often problems in the spleen.
When the heart is under strain, the concerned person can react overly emotional, even hysterical.
Also in this case, the middle of the back is concerned, especially on the height of the heart.

There's also another view on the meridians\index{meridian} and the organs, which is way more abstract and spiritual, but as valid and clear.
To understand this type of diagnosis, we have to first recognize, that every organ has a physical function, but also an abstract function.
The work abstract is here meaning the spiritual dimension.

A simple example: the job of our digestive system is to take in the nutrition from the exterior world, digest it, make the nutrients accessible to the blood and eventually excrete the unnecessary parts.
That is in it's own right a spiritual task!
In the basic principle, this means the capacity of a human being to take what it needs, use of it what is necessary for survival and happiness and to excrete the superfluous part.

This process has a big signification for our physical, mental and spiritual well being.
Maybe you ask yourself now: 
"Do I receive what I need for my life? And how am I able to get rid of unnecessary things, like experiences and habits, which don't correspond to my current level of development anymore?"
When you answered these questions, think about how good your digestion really is.

So the digestive system is really just the physical expression of a spiritual function. 
Every organ can be looked at this way.
All organs and bodily functions, which compose your body are really just physical manifestations of spiritual qualities.
These spiritual qualities were already manifest in your soul by the time it connected to earth during your conception.

Let's look at a second example, the kidneys.
The main function of the kidneys are to filter impurities from the blood.
The kidneys are part of the vital organs.
In an abstract sense, the cleaning function of the kidneys is based on it recognizing what is good and what is unnecessary or even harmful.
If the kidneys don't work properly we have a very difficult time to distinguish between good and bad in life.

Put in different words: The discernment will be impaired if the kidneys don't work properly. 
Old unsavory things about ourselves, our personality and surrounding are conserved, even toxic elements.
We can't decide, if we should admit to it or put them aside.
As a consequence is that we constantly are going to be surprised by unforeseen difficulties, sometimes even shocked.
That leads to anxiety, in the extreme case even to delusions.
We feel as victim of life, whereas we are in reality victims of our lacking discernment.

In each organ and meridian there's the risk, that too much energy accumulates.
On the other hand, it can also lack in energy, that's what the Japanese work Kyo describes.\index{Kyo}
If an organ has too much energy, then the activity and meridian of this organ is in turn excessively strong.
The organ can get overloaded.

In addition to that, when the energy gets blocked in this place, there's a congestion created.
Then other organs are not sufficiently supplied with energy, even though there would be enough energy available.
The mental/psychological aspect of this overactive organ will be in turn overemphasized.
Is an organ Kyo or exhausted, it will get weak and sluggish, the blood will congest there.
The mental aspect shows that we are weakened in this part of our lives, which corresponds to the weakened organ.

If I know for instance the symptoms of a weak stomach, then I can in this case ask a friend suffering from it if they suffer from the mental experiences, which are associated with this organ.
Then i can recommend to strengthen the corresponding organ, like in this case the stomach, which in turn will also cure the mental well being.

So by learning about the abstract nature of the organs and meridians you can learn a lot about the mental and psychological well being a a person.

In the following chapter, the meridians and their corresponding physical aspects are listed and explained.
The mostly mental aspects are represented in a circle.
\end{document}